% Volume IX: Institutional Design
% Part XVII. Designing for Delayed Truth
% Chapter 100: Attention-Aware Archives
% Spine: proof-spines/ch100-spine.md

\chapter{Attention-Aware Archives}
\label{ch:ch100}

An archive fit for delayed truth must be designed for later recovery, not merely present storage. It should expose provenance graphs, uncertainty annotations, duplicate clusters, dissenting classifications, and discovery histories so that a later investigator can see how a record traveled, what was contested, what may have been overcounted, and what remained obscure.

Search is part of that architecture. If earlier salience models suppressed an item, the archive should still make it recoverable rather than ratifying the first ranking as destiny. Attention-aware archives therefore treat findability as a condition of accountability: what matters is not only whether information exists somewhere, but whether future inquiry can locate it despite the biases, omissions, and heuristics of earlier attention.
